% Acuerdos incorporados; integrar al catálogo AMFE original cuando esté disponible.
\begin{formulario}{T-16}

Los registros siguientes fijan las condiciones acordadas de capacitación y
adhesión. Los códigos AC-CAP y AC-ADH identifican estos acuerdos; no reemplazan
ni renumeran riesgos de un catálogo previo. \marcador{Integrar al catálogo de
riesgos original y verificar su cuantificación; fuente integrada no disponible}.

\begin{formularioTDieciseis}
\riesgoTDieciseis{
 id = {AC-CAP},
 riesgo = {A13 no puede iniciar por falta de versión de entrenamiento y material probado y congelado en PREPROD},
 categoria = {Proyecto},
 probabilidad = {\marcador{Por cuantificar}},
 impacto = {\marcador{Por cuantificar}},
 exposicion = {\marcador{Por cuantificar}},
 mitigacion = {G13 al cierre del sprint 7; funciones certificadas definidas con D3, datos ficticios y material por perfil. A13: A06, A09 FS$-10$, G13; 40 días, diez de familiarización y treinta de práctica y evaluación, incluido complemento del sprint 8 antes de marcha blanca. No iniciar sin G13; registrar atraso y recalcular riesgo con la misma red del T-15},
 responsable = {Jefe de Proyecto y Líder de Implantación}
}
\riesgoTDieciseis{
 id = {AC-ADH},
 riesgo = {59 o menos transportistas firmados de la población inicial de 148 al cierre de M9},
 categoria = {Contraparte},
 probabilidad = {\marcador{Por cuantificar}},
 impacto = {\marcador{Por cuantificar}},
 exposicion = {\marcador{Por cuantificar}},
 mitigacion = {Corte: 31 de octubre de 2027 bajo SA-01. Una vez por RUT; anexo firmado por mandante y representante autorizado, vigente, no revocado y registrado al corte. Ambas modalidades. Refuerzo con menos de 60: ampliar reparto de espera, priorizar retornos y extender modalidad de datos. Metas: 104 en M16 y 134 en M21. Instalación, activación y consentimientos se controlan por separado; la firma del transportista no sustituye los consentimientos de conductores},
 responsable = {Coordinación de adhesión y mandante}
}
\end{formularioTDieciseis}
\end{formulario}
